\section{AI Infra 架构思路（本机全量优先）}

\begin{designbox}{定位校准}
AI Infra 是 SciScope 的工程可靠性和可复现性保障。国赛演示以\textbf{本机全量}为主路径：
本地 PostgreSQL/pgvector 承载 16 万论文与 36 万+ 片段, 本地 embedding/reranker 提供语义能力,
生成模型走 DeepSeek API 或本地 LLM, 全部通过 Makefile 一键可复现。
\end{designbox}

\subsection{运行 profile（三档）}

\begin{center}
\begin{tabularx}{\textwidth}{>{\bfseries}p{28mm}X X}
\toprule
Profile & 用途 & 技术路径 \\
\midrule
Full（主演示） & 国赛本机全量：检索/核查/争议/MCP & 本地 PostgreSQL + pgvector + E5 + reranker + DeepSeek \\
Lite & 无模型快速启动与 CI 验收 & Mock provider + 预置证据 + 样本语料 \\
Hosted demo & 公开入口（小样本） & Render + Supabase + DeepSeek, 关闭运行时 embedding \\
\bottomrule
\end{tabularx}
\end{center}

\begin{plainbox}
\textbf{口径纪律}: hosted demo 约 220 篇/467 片段, 不等于国赛演示能力全集。国赛现场以本机全量为准;
提交配套"评委视频 + 硬件清单"双轨, 说明全量复现的门槛与关键命令, 不拿 hosted 冒充全量。
\end{plainbox}

\subsection{服务层拓扑}

\begin{tikzpicture}[
  node distance=6mm,
  every node/.style={font=\small},
  box/.style={draw=black, fill=white, rounded corners=2pt, align=center, minimum width=36mm, minimum height=11mm},
  arrow/.style={-{Latex[length=2mm]}, thick, color=black}
]
  \node[box, very thick] (pg) {PostgreSQL + pgvector\\papers / chunks / embeddings / stance};
  \node[box, right=10mm of pg] (fastapi) {FastAPI 后端\\检索 / 证据 / Agent / MCP};
  \node[box, below=12mm of fastapi] (tui) {Go TUI（SSE 消费）};
  \node[box, left=14mm of fastapi, densely dashed] (llm) {生成层\\DeepSeek API / 本地 LLM / Mock};
  \draw[arrow] (pg) -- (fastapi);
  \draw[arrow] (fastapi) -- (tui);
  \draw[arrow,densely dashed] (llm) -- (fastapi);
\end{tikzpicture}

\begin{plainbox}
\begin{itemize}
  \item PostgreSQL/pgvector 是\textbf{正式服务层}, 不是可选项; 本地全量检索与 stance 落库都依赖它。
  \item 本地 E5 与 bge-reranker 是语义臂与重排器; 不可用时检索自动降级, 核查回退仅判"证据不足"。
  \item Makefile 是唯一工程入口: \texttt{make postgres-refresh}, \texttt{make agent-build},
        \texttt{make backend}, \texttt{make tui}, \texttt{make eval-stance}, \texttt{make submission-package}。
  \item 会话状态为进程内 MemorySaver + JSON 落盘（data/.agent\_memory）, 不做数据库级持久化——如实声明。
\end{itemize}
\end{plainbox}

\subsection{模型 provider 路由}

\begin{center}
\begin{tabularx}{\textwidth}{p{33mm}X X}
\toprule
Provider & 适合场景 & 注意事项 \\
\midrule
Mock & CI、无网、接口验收 & 必须明确标记, 不代表最终回答质量 \\
DeepSeek API & 正式生成、中文综述、复杂推理 & 需要 API Key 与网络环境 \\
本地 LLM（Qwen 7B 4bit） & 本机离线演示（可选下载） & 需要模型下载与预热 \\
\bottomrule
\end{tabularx}
\end{center}

\subsection{Infra 应该服务的产品指标}

\begin{plainbox}
\begin{itemize}
  \item \textbf{启动确定性}: 一条命令启动对应 profile, 端口和模型状态可检查（\texttt{make tui-doctor}）。
  \item \textbf{回答可追溯}: 无论 provider 是谁, 答案必须绑定本地检索证据。
  \item \textbf{失败可降级}: DeepSeek 不可用时切本地 LLM 或 Mock; 数据库不可用时 stance 写入静默跳过。
  \item \textbf{诚实缺口}: 无鉴权/租户、无 OTel、会话非多实例持久化——原型级治理, 报告中如实声明。
\end{itemize}
\end{plainbox}

\clearpage